% SPDX-License-Identifier: Apache-2.0
%
% The repository in numbers: lines per topic, measured once and typed
% in. Built by //docs:stats.
\documentclass[journal,onecolumn,9pt]{IEEEtran}

% One column: 80 columns fit at the body size, so nothing is reduced.
\newcommand{\listingsize}{\normalsize}
% SPDX-License-Identifier: Apache-2.0
%
% The house preamble, shared by every document in this directory. Loaded
% after \documentclass, before \title. Keeping it in one file is what
% stops two articles drifting apart in font, listing style or figure
% style.
% Computer Modern. IEEEtran selects Times (ptm) by default, so the face has
% to be chosen explicitly.
%
% Latin Modern is the Computer Modern variety used here, and the choice is
% forced rather than aesthetic. Under T1 encoding, plain `cmr` has no Type 1
% outlines unless cm-super is installed, and it is not in the pinned
% distribution, so pdflatex falls back to EC bitmaps and every glyph in the
% PDF becomes a Type 3 font. That was measured: the first build produced a
% document with 10 Type 3 fonts and no Type 1. Latin Modern is Computer
% Modern redrawn with a full T1 repertoire and real .pfb outlines, so the
% same design comes out scalable.
\usepackage[T1]{fontenc}
\usepackage[utf8]{inputenc}
\usepackage{lmodern}
% microtype is not in the pinned TeX distribution. emergencystretch alone
% keeps the two-column measure from producing overfull lines.
\setlength{\emergencystretch}{3em}

\usepackage{amsmath}
\usepackage{amssymb}
\usepackage{amsthm}
\usepackage{booktabs}
\usepackage{array}
% enumitem is not in the pinned TeX distribution; the list spacing below
% does what \setlist would have done.
\usepackage{float}
\usepackage{xcolor}
\usepackage{listings}
\usepackage{tikz}
\usetikzlibrary{arrows.meta,positioning,fit,backgrounds,calc}
\usepackage[hidelinks,bookmarks=true]{hyperref}

% Numbered, definition-style box for a rule the reader is meant to apply.
\theoremstyle{definition}
\newtheorem{rulebox}{Rule}
\newtheorem{example}{Example}

% Unnumbered asides.
\theoremstyle{remark}
\newtheorem*{tip}{Tip}
\newtheorem*{pitfall}{Pitfall}

% Tighter lists, in place of enumitem's \setlist.
\setlength{\topsep}{2pt}
\setlength{\itemsep}{1pt}
\setlength{\parsep}{0pt}

\newcommand{\code}[1]{\texttt{#1}}
\newcommand{\term}[1]{\emph{#1}}

% Syntax highlighting. Four roles, distinguishable in colour and still
% distinguishable in greyscale, because an IEEE article gets printed: the
% keyword is the only bold role, the comment is the only italic one, and the
% two remaining roles differ in lightness as well as in hue.
\definecolor{kwblue}{RGB}{20,60,150}
\definecolor{cmtgreen}{RGB}{90,120,90}
\definecolor{strred}{RGB}{150,50,40}
\definecolor{numgray}{gray}{0.45}
\definecolor{framegray}{gray}{0.70}
\definecolor{bgshade}{gray}{0.97}
% A darker fill for figure cells. The listing background has to stay very
% light to keep code readable; a figure cell has to be clearly shaded or the
% distinction it draws is invisible in print.
\definecolor{cellshade}{gray}{0.90}

% The listing size is the document's choice, because it depends on the
% column: a two-column article sets `\listingsize` to 6pt and keeps its
% listings within 70 columns, a one-column document keeps the body size
% and includes source files at 80. Lines are never broken; a line that does not fit
% overflows the frame, which is visible, rather than wrapping, which is
% not.
\providecommand{\listingsize}{\scriptsize}

% `'` is deliberately not a string delimiter: the language uses it for loop
% labels, as Rust does, so treating it as a quote would swallow the rest of
% the line.
\lstdefinestyle{txhdl}{
  basicstyle=\ttfamily\listingsize,
  keywordstyle=\bfseries\color{kwblue},
  commentstyle=\itshape\color{cmtgreen},
  stringstyle=\color{strred},
  numberstyle=\tiny\color{numgray},
  morekeywords={mod,use,pub,const,type,struct,enum,trait,impl,fn,async,let,mut,
    match,if,else,loop,while,for,in,break,continue,return,as,await,
    module,bus,channel,transaction,protocol,pipeline,flow,seq,config,
    port,stage,pipe,instance,connect,bind,tag,domain,blackbox,foreign,
    property,role,signal,handler,true,false,
    sequence,interface,modport,implements,package,import,var,wire,func,proc,
    case,default,next,fork,branch,join,state,fsm},
  morecomment=[l]{//},
  morestring=[b]",
  showstringspaces=false,
  columns=fullflexible,
  keepspaces=true,
  breaklines=false,
  % Framed, so a listing is bounded rather than floating in the column.
  frame=single,
  rulecolor=\color{framegray},
  framesep=3pt,
  backgroundcolor=\color{bgshade},
  xleftmargin=3pt,
  xrightmargin=3pt,
  aboveskip=6pt,
  belowskip=6pt,
  % Every listing is numbered and captioned, so it can be referred to by
  % number. `caption` without `float` numbers the listing and leaves it
  % where it was written, which is what a listing in running prose needs.
  captionpos=b,
  abovecaptionskip=2pt,
  belowcaptionskip=6pt,
}

% Figures drawn as text. Framed the same way, and with the highlighting
% switched off, because a box-and-arrow diagram is not code and half its
% words turning blue reads as an accident. Setting a style adds keys rather
% than clearing the ones already set, so the three roles are neutralised by
% hand rather than by leaving them out.
\lstdefinestyle{ascii}{
  basicstyle=\ttfamily\listingsize,
  keywordstyle=\ttfamily,
  commentstyle=\ttfamily,
  stringstyle=\ttfamily,
  columns=fullflexible,
  keepspaces=true,
  frame=single,
  rulecolor=\color{framegray},
  framesep=3pt,
  backgroundcolor=\color{bgshade},
  xleftmargin=3pt,
  xrightmargin=3pt,
  aboveskip=6pt,
  belowskip=6pt,
}
\lstset{style=txhdl}

% House TikZ styles. `shorten` lives on the arrow styles rather than on each
% arrow, so every arrow in the document gets the gap, including ones added
% later. TikZ clips a path at a node's border, so without it an arrowhead
% butts against the box with no gap at all. Keep `node distance` well above
% twice the shorten, or the arrow shrinks to nothing.
\tikzset{
  box/.style={draw,rounded corners=2pt,align=center,inner sep=4pt,
              font=\footnotesize},
  boxf/.style={box,fill=bgshade},
  lbl/.style={font=\scriptsize\itshape,align=center},
  fl/.style={-{Stealth[length=4pt]},shorten >=2.5pt,shorten <=2.5pt},
  fld/.style={fl,dashed},
}



\InputIfFileExists{buildstamp.tex}{}{}
\providecommand{\buildstamp}{Draft build (outside the Bazel flow).}

\title{TxHDL by the Numbers}

\author{Filip Filmar and Dragiša Janković%
\thanks{Both authors are independent. Filip Filmar is the
corresponding author, at f@hdlfactory.com; Dragiša Janković is
reached at linkedin.com/in/dragisa-jankovic.}%
\thanks{This document is a snapshot. The numbers in it were measured
on September 13, 2026, at commit \code{0ab36f3} of the mainline, and
typed in; nothing rebuilds them, and they are not kept up to date.
The script that produced them is in the tree, for a reader who wants
the numbers of a later day.}%
\thanks{The authors used a large language model, Claude, as an
assistant in exploring the concepts and in writing and constructing
the documents and the programs; every commit of the repository says
so and carries the prompts that led to it, verbatim.}%
\thanks{\buildstamp}}

\markboth{TxHDL by the Numbers}{Filmar and Janković: TxHDL by the Numbers}

\begin{document}
\maketitle

\section{What Was Counted}

Every file tracked by git in the repository \code{HDL/txhdl}, outside
\code{third\_party/}, was counted with \code{wc -l}, so comments and
blank lines count as lines. The lock files, \code{Cargo.lock},
\code{MODULE.bazel.lock}, \code{cargo-bazel-lock.json} and
\code{multitool.lock.json}, and \code{.bazelversion} were left out:
they are generated, and they would have added some thirty thousand
lines that nobody wrote. \code{third\_party/} holds vendored \LaTeX{}
packages and the IEEE document class, under their own licenses, and is
not this project's work.

The measurement was taken on September 13, 2026, at commit
\code{0ab36f3} of the mainline, the commit before this document and
its script were added. Listing~\ref{lst:stats-script} is the script;
run from the root of a checkout, it prints
Tables~\ref{tab:stats-topic} and~\ref{tab:stats-lang} for that
checkout. The topics are the directories of the tree, named in the
script.

\lstinputlisting[style=ascii,caption={\code{tools/repostats.sh}: the
measurement.},label={lst:stats-script}]{tools/repostats.sh}

\section{Lines per Topic}

Table~\ref{tab:stats-topic} is the count by topic, largest first.
The total is 218 files and 32,452 lines.

\begin{table}[h]
\caption{Lines per topic, September 13, 2026.}
\label{tab:stats-topic}
\centering
\begin{tabular}{@{}llrr@{}}
\toprule
Topic & Where & Files & Lines \\
\midrule
Documents, \LaTeX{}
  & \code{docs/*.tex}, \code{docs/*\_sections/}
  & 74
  & 7,109 \\
Early specs and proposals
  & \code{spec/}, \code{filmil/}, \code{proposal.md}
  & 8
  & 5,237 \\
Runtime library & \code{lib/src/} & 7 & 3,733 \\
Vreteno core & \code{cpu/vreteno/src/} & 12 & 2,951 \\
Examples & \code{lib/examples/} & 32 & 2,878 \\
Lowering macro & \code{lib/macros/lib.rs} & 1 & 2,556 \\
Tools & \code{tools/} & 16 & 2,091 \\
Design notes, Markdown & \code{docs/*.md} & 5 & 1,468 \\
Docs build files
  & \code{docs/BUILD.bazel}, \code{docs/waveform.bzl}
  & 2
  & 1,162 \\
Rust probes & \code{experiments/} & 34 & 847 \\
Root config and READMEs & the root & 9 & 693 \\
Vreteno board & \code{cpu/vreteno/board/} & 7 & 504 \\
Lib build files
  & \code{lib/BUILD.bazel}, \code{lib/foreign.bzl} and two more
  & 5
  & 408 \\
Vreteno lockstep test & \code{cpu/vreteno/tests/} & 1 & 313 \\
CI workflows & \code{.forgejo/workflows/} & 3 & 306 \\
Vreteno build and synthesis
  & \code{cpu/vreteno/BUILD.bazel}, \code{vreteno.xdc}
  & 2
  & 196 \\
\midrule
Total & & 218 & 32,452 \\
\bottomrule
\end{tabular}
\end{table}

\section{Lines per Language}

Table~\ref{tab:stats-lang} is the same count by file extension.

\begin{table}[h]
\caption{Lines per language, September 13, 2026.}
\label{tab:stats-lang}
\centering
\begin{tabular}{@{}llr@{}}
\toprule
Extension & Language & Lines \\
\midrule
\code{rs} & Rust & 14,893 \\
\code{md} & Markdown & 7,293 \\
\code{tex} & \LaTeX{} & 7,109 \\
\code{bazel} & Bazel, \code{BUILD.bazel} and \code{MODULE.bazel} & 2,022 \\
\code{bzl} & Starlark & 317 \\
\code{yml} & workflow YAML & 306 \\
\code{sh} & shell & 167 \\
\code{v} & Verilog & 148 \\
\code{go} & Go & 119 \\
\code{xdc} & Vivado constraints & 33 \\
\code{vhdl} & VHDL & 20 \\
other
  & \code{.bazelrc}, \code{main.cc}, \code{rustfmt.toml}, \code{.gitignore}
  & 25 \\
\bottomrule
\end{tabular}
\end{table}

\section{The Largest Files}

Table~\ref{tab:stats-files} is the runtime and the core file by file,
with the tools, since those are the files a reader asks about. The
largest file in the tree is the lowering macro; the second is the
early language specification, which predates the library and is kept
as history.

\begin{table}[h]
\caption{The runtime, the core and the tools, file by file.}
\label{tab:stats-files}
\centering
\begin{tabular}{@{}lr@{\hspace{12mm}}lr@{\hspace{12mm}}lr@{}}
\toprule
Runtime & Lines & Vreteno & Lines & Tools & Lines \\
\midrule
\code{macros/lib.rs} & 2,556 & \code{core.rs} & 815 & \code{vshim} & 579 \\
\code{comp.rs} & 1,697 & \code{isa.rs} & 522 & \code{fst2tb} & 547 \\
\code{netlist.rs} & 1,169 & \code{model.rs} & 364 & \code{dt2tikz} & 372 \\
\code{types.rs} & 479 & \code{program.rs} & 355 & \code{api} & 259 \\
\code{foreign.rs} & 200 & \code{uart.rs} & 246 & \code{fstcheck} & 170 \\
\code{pipeline.rs} & 95 & \code{main.rs} & 221 & & \\
\code{funcs.rs} & 48 & \code{term.rs} & 100 & & \\
\code{lib.rs} & 45 & \code{timer.rs} & 98 & & \\
 & & \code{router.rs} & 98 & & \\
 & & \code{dmem.rs} & 78 & & \\
 & & \code{bus.rs} & 41 & & \\
 & & \code{lib.rs} & 13 & & \\
\bottomrule
\end{tabular}
\end{table}

\section{What the Numbers Say}

The language is small. The runtime and the macro together are about
6,300 lines of Rust, and of the runtime, \code{comp.rs} and
\code{netlist.rs} are three quarters; the rest is the types, the
foreign units, the pipeline and the functions.

Vreteno is about half the size of the language. The core proper is
815 lines, the decoder 522 and the reference model 364; the devices on
the bus, the UART, the timer, the router and the data memory, are
about 560 together.

Prose outweighs code. The documents, the design notes and the early
specifications sum to about 13,800 lines, against about 15,000 of
Rust, and the documents include their listings from the tree rather
than repeating them, so the prose is prose.

The examples are many and short: thirty-two files of ninety lines on
average, one per concept, which is the rule of the project.

The tools are a quarter of the Rust, and the two testbench generators,
\code{vshim} and \code{fst2tb}, are the largest of them at about 550
lines each: the checks are as big a part of the work as the checked.

\end{document}
